\section{Visual notation}

% ── slide p.1 ──────────────────────────────────────────
% SKIP: title_slide

% ── slide p.2 ──────────────────────────────────────────
\subsection{Objective}

A visual notation is the interface between domain experts, who own the processes, and formal analysis tools, which need an unambiguous representation. The reference for the chapter is Ch.~3 of Weske.\footnote{Weske M., \emph{Business Process Management: Concepts, Languages, Architectures}, Springer, 2007, Ch.~3.}

% ── slide p.3 ──────────────────────────────────────────
Section~2 collected three literature definitions of business process (Hammer \& Champy, Davenport, Johansson). The textbook definition below distills them and is the one used throughout these notes:

\begin{definitionbox}{Business process}
A \textbf{business process} consists of a set of activities performed in coordination within an organisational and technical environment. These activities jointly realise a business goal.\footnote{Weske M., \emph{Business Process Management}, Springer, 2007.}
\end{definitionbox}

Each business process is enacted by a single organisation, but it may interact with processes performed by other organisations. The intra-organisational view is called \emph{orchestration}; the inter-organisational one is \emph{collaboration} (or \emph{choreography}).

% ── slide p.4 ──────────────────────────────────────────
\begin{definitionbox}{Business process model and instance}
A \textbf{business process model} consists of a set of activity models together with the execution constraints between them. A \textbf{business process instance} represents one concrete case in the operational business of a company, consisting of activity instances.\footnote{Weske M., \emph{Business Process Management}, Springer, 2007.}
\end{definitionbox}

% ── slide p.5 ──────────────────────────────────────────
An activity model is a \emph{blueprint} for its activity instances, exactly as a process model is a blueprint for its process instances, one per \emph{case}: each instance carrying its own case data, such as customer identifier, name, and amount of the credit request. When no confusion is possible, \textbf{activity} denotes both an activity model and an activity instance, and \textbf{process} likewise covers both sides of the model/instance distinction.

% ── slides p.6--p.10 ────────────────────────────────────
\subsection{What a model must capture}

Modelling starts from five questions about the process.

\begin{itemize}
  \item \textbf{Who is the customer?} Every process starts and ends with a customer, who requests a product and ultimately receives it. The customer may be another department of the same company: the role is independent of organisational boundaries.
  \item \textbf{Who is the owner?} The \emph{process owner} is responsible for the process: that instances run correctly, that business goals are met, that performance is measured and improved. The model has to support these responsibilities, not contradict them.
  \item \textbf{Which tasks?} The \emph{activities} (or \emph{tasks}) needed to realise the goals. A task is treated as an atomic action, possibly with a duration and a cost, but granularity is a modelling choice: any task may decompose into finer steps on closer inspection.
  \item \textbf{Which dependencies?} \emph{Execution constraints} fix the admissible orderings: which task precedes which, which run in parallel, which exclude each other. They distribute work over \emph{time}; a \emph{process orchestration language} is the formalism they are written in.
  \item \textbf{Which roles?} Each task may require specific abilities (a customer-service agent, a back-office clerk, a software service) abstracted as \emph{roles}. Roles distribute work over \emph{space}: which actor is in charge of which step. Time and space are orthogonal dimensions of one orchestration.
\end{itemize}

% ── slide p.11 ──────────────────────────────────────────
Modelling translates informal textual requirements into a notation, ideally graphical. An explicit diagram lets stakeholders \textbf{communicate} about how work is done, \textbf{refine} the process when ambiguities surface, and \textbf{improve} it over time.

% ── slide p.12 ──────────────────────────────────────────
\begin{definitionbox}{Business process management}
\textbf{Business process management} (BPM) is the discipline that brings together concepts, methods, and techniques to support the design, administration, configuration, enactment, and analysis of business processes.\footnote{Weske M., \emph{Business Process Management}, Springer, 2007.}
\end{definitionbox}

BPM rests on an \emph{explicit representation} of processes, tasks, and execution constraints. Once the representation is explicit, the process can be analysed (does it terminate? can it deadlock?), improved (where are the bottlenecks?), and enacted by software.

% ── slide p.13 ──────────────────────────────────────────
\begin{definitionbox}{Business process management system}
A \textbf{business process management system} (BPMS) is a generic software system, driven by explicit process representations, that coordinates the enactment of business processes.\footnote{Weske M., \emph{Business Process Management}, Springer, 2007.}
\end{definitionbox}

Process models are the main artefact for implementing a process. The implementation may rest on organisational rules and written procedures alone, or be delegated to a BPMS that reads the model and steers the instances through it: assigning tasks to roles, enforcing execution constraints, logging what actually happened.

% ── slide p.14 ──────────────────────────────────────────
\subsection{Representing processes}

Process representations come in three flavours, one per audience:

\begin{itemize}
  \item \textbf{Visual representations}: diagrams for \emph{humans}, informal but intuitive: BPMN, EPC, graphical fragments of BPEL. What we see.
  \item \textbf{Languages}: unambiguous machine syntax, process dialects and XML schemas such as BPEL or BPMN~XML. What machines see.
  \item \textbf{Models}: rigorous semantics for \emph{analysis}, automata, Petri nets, workflow nets. What we analyse.
\end{itemize}

The layers complement each other, and a well-engineered BPM stack provides translations between them.

% ── slide p.15 ──────────────────────────────────────────
XML, the \emph{eXtensible Markup Language}, is the serialisation format most process languages adopt. An XML document is a tree of user-defined \emph{tags} with the values at the leaves:

\begin{verbatim}
<?xml version="1.0" encoding="UTF-8"?>
<note>
  <from>Alice</from>
  <to>Bob</to>
  <heading>Reminder</heading>
  <body>Don't forget to buy oranges!</body>
</note>
\end{verbatim}

Verbose but unambiguous, and readable by any standard parser.

% ── slides p.12--p.14 (Weske BPM lifecycle) ─────────────
\subsection{The BPM lifecycle}

The BPM definition lists five capabilities: design, administration, configuration, enactment, analysis. These are not a checklist run once but a \emph{cycle} that feeds its measurements back into the next redesign, closing back onto its own start. Each turn is a phase with a characteristic activity, a course technique that lives there, and a responsible role. Roles matter especially at \emph{validation}: a model is only as good as the agreement between the \emph{owner} who states the goals, the \emph{analyst} who draws the diagram, and the \emph{participants} who will run it: which is why the design phase closes with the three arguing over one shared diagram instead of over prose requirements.
\begin{itemize}
  \item \textbf{Design \& Analysis.} The process is modelled and checked before anyone runs it. Modelling produces a diagram in a graphical notation, \textbf{BPMN} or \textbf{EPC}; validation confronts it with stakeholders and, informally, with simulation of sample cases; verification translates it into a \textbf{workflow net} and proves properties such as \emph{soundness}. Role: the \emph{process analyst}, with the \emph{owner} who states the goals and the \emph{participants} who confirm the model matches the real work.
  \item \textbf{Configuration.} The validated model becomes an executable process: a \textbf{BPMS} is selected and configured, each task bound to the roles and software services that perform it, the model made executable by a compliant engine (the \texttt{.bpmn} serialisation is the exchange format). Role: the \emph{system architect}, wiring the model into the platform.
  \item \textbf{Enactment.} The configured process runs: the BPMS creates one instance per case, assigns tasks, enforces the execution constraints, and \emph{logs} every step that actually happens. Role: the \emph{participants} who carry out the tasks, steered and recorded by the BPMS.
  \item \textbf{Evaluation.} The logs feed back into analysis. \textbf{Process mining} reconstructs the real process from the log; \textbf{conformance} checking compares it against the designed model; \emph{quantitative analysis} measures durations, costs, bottlenecks. The findings drive the next redesign, returning to Design \& Analysis. Role: the \emph{process analyst}, reporting to the \emph{owner} who decides what to improve.
\end{itemize}
Design and Evaluation are the two analytic phases, and they use different evidence: Design \& Analysis reasons on the \emph{model} before execution (a workflow net checked for soundness), Evaluation on the \emph{log} after execution (a real trace, mined and compared). The former asks whether the process \emph{can} misbehave; the latter, whether it \emph{did}.

% ── slide p.16 ──────────────────────────────────────────
\subsection{BPMN vs .bpmn}

The visual/machine contrast is concrete in BPMN. A one-task process: start event, activity \emph{Run}, end event:

\begin{center}
\begin{tikzpicture}[node distance=11mm]
  \node[bpmn event] (s) {};
  \node[bpmn task, right=of s] (run) {Run};
  \node[bpmn event, line width=1.1pt, right=of run] (e) {};
  \draw[bpmn flow] (s) -- (run);
  \draw[bpmn flow] (run) -- (e);
\end{tikzpicture}
\end{center}

Serialised as a \texttt{.bpmn} file, each visual element becomes an XML tag with a unique identifier (\verb|StartEvent_1|, \verb|Activity_1lzhm01|, \verb|Event_1t0u7im|) and each sequence flow becomes a \verb|<bpmn:sequenceFlow>| tag whose \verb|sourceRef| and \verb|targetRef| attributes point at those identifiers. The two representations are interchangeable: the editor renders the XML as the diagram, and saving regenerates it.

% ── slides p.17--p.18 ───────────────────────────────────
\subsection{Insurance claim example}

The running example for the rest of the section is a simplified insurance-claim process.\footnote{van der Aalst W., van Hee K., \emph{Workflow Management: Models, Methods, and Systems}, MIT Press, 2004, Sect.~1.3.} Its textual description lists twelve steps:

\begin{enumerate}
  \item \textbf{recording} the receipt of the claim;
  \item establishing the \textbf{type} of the claim;
  \item checking the covering of the client's \textbf{policy};
  \item checking the \textbf{premium} (are payments up to date?);
  \item decision for \textbf{rejection}/admission;
  \item if step 3 or step 4 has a negative result: producing a \textbf{rejection letter}, then jump to step 12;
  \item if steps 3 and 4 have positive results: sending an \textbf{estimate} of the amount to be paid;
  \item recording the client's \textbf{reaction};
  \item \textbf{assessment} of an objection;
  \item if step 9 has a negative result: decision to \textbf{revise} step 7;
  \item if step 9 has a positive result: \textbf{payment} of the claim;
  \item \textbf{filing} and closure of the claim.
\end{enumerate}

The description mixes activities (recording, filing, payment) with control flow (``if/then'', ``then 12'') and feedback (revise step 7).

% ── slides p.19--p.20 ───────────────────────────────────
\subsection{Tasks and links}

Each step becomes a task node; two of them, \textsc{rejection?} and \textsc{assessment?}, are decision points whose outcome is a boolean. Links then wire the tasks into a flow: a sequence, a parallel section, two choices, and a feedback loop:

\begin{center}
\begin{tikzpicture}
  \node[bpmn task, font=\scriptsize\sffamily] (rec) at (0,0) {1.~recording};
  \node[bpmn task, font=\scriptsize\sffamily] (typ) at (2.1,0) {2.~type};
  \node[bpmn task, font=\scriptsize\sffamily] (pol) at (4.2,0.9) {3.~policy};
  \node[bpmn task, font=\scriptsize\sffamily] (pre) at (4.2,-0.9) {4.~premium};
  \node[bpmn task, font=\scriptsize\sffamily] (rej) at (6.4,0) {5.~rejection?};
  \node[bpmn task, font=\scriptsize\sffamily] (let) at (8.9,1.1) {6.~reject letter};
  \node[bpmn task, font=\scriptsize\sffamily] (est) at (8.9,-1.1) {7.~estimate};
  \node[bpmn task, font=\scriptsize\sffamily] (rea) at (8.9,-2.4) {8.~reaction};
  \node[bpmn task, font=\scriptsize\sffamily] (ass) at (11.3,-2.4) {9.~assessment?};
  \node[bpmn task, font=\scriptsize\sffamily] (rev) at (11.3,-1.1) {10.~revision};
  \node[bpmn task, font=\scriptsize\sffamily] (pay) at (13.7,-2.4) {11.~payment};
  \node[bpmn task, font=\scriptsize\sffamily] (fil) at (13.7,1.1) {12.~filing};
  \draw[bpmn flow] (rec) -- (typ);
  \draw[bpmn flow] (typ) -- (pol);
  \draw[bpmn flow] (typ) -- (pre);
  \draw[bpmn flow] (pol) -- (rej);
  \draw[bpmn flow] (pre) -- (rej);
  \draw[bpmn flow] (rej) -- (let);
  \draw[bpmn flow] (rej) -- (est);
  \draw[bpmn flow] (let) -- (fil);
  \draw[bpmn flow] (est) -- (rea);
  \draw[bpmn flow] (rea) -- (ass);
  \draw[bpmn flow] (ass) -- (rev);
  \draw[bpmn flow] (rev) -- (est);
  \draw[bpmn flow] (ass) -- (pay);
  \draw[bpmn flow] (pay) -- (fil);
\end{tikzpicture}
\end{center}

% ── slides p.21--p.26 ───────────────────────────────────
\subsection{Control-flow patterns}

The insurance net combines four recurring control-flow patterns, visible as fragments of the diagram above.

\begin{itemize}
  \item \textbf{Sequence}: \textsc{type} can start only after \textsc{recording} has finished; a single arrow makes the temporal precedence explicit.
  \item \textbf{Parallel}: after \textsc{type}, both \textsc{policy} and \textsc{premium} must be checked before \textsc{rejection?} can be evaluated. The two branches run independently (their relative order is irrelevant) and the join waits for \emph{both} (AND-split / AND-join).
  \item \textbf{Choice}: from \textsc{rejection?} the flow goes \emph{either} to \textsc{reject letter} \emph{or} to \textsc{estimate}, depending on the boolean outcome; exactly one branch fires per instance (XOR-split). The same construct appears at \textsc{assessment?}.
  \item \textbf{Merge}: the two alternative paths reunite at \textsc{filing}, reached through \textsc{reject letter} or through \textsc{payment}. Unlike the parallel join, the merge does not wait for both branches: whichever arrives carries the flow through (XOR-join).
  \item \textbf{Iteration}: a negative \textsc{assessment?} routes back through \textsc{revision} to \textsc{estimate}; the cycle \textsc{estimate} $\rightarrow$ \textsc{reaction} $\rightarrow$ \textsc{assessment?} $\rightarrow$ \textsc{revision} repeats until the objection is finally accepted or rejected. Iteration combines an XOR-join (entering the loop body) with an XOR-split (repeat or exit).
\end{itemize}

% ── slides p.27--p.29 ───────────────────────────────────
\subsection{Ambiguity}

The diagram uses one kind of arrow for two different forks, and this is genuinely ambiguous. After \textsc{type} both branches are always executed; after \textsc{rejection?} exactly one is; yet the two fan-outs look identical:

\begin{center}
\begin{tikzpicture}
  \node[bpmn task, font=\scriptsize\sffamily] (t) at (0,0) {2.~type};
  \node[bpmn task, font=\scriptsize\sffamily] (p3) at (2.6,0.7) {3.~policy};
  \node[bpmn task, font=\scriptsize\sffamily] (p4) at (2.6,-0.7) {4.~premium};
  \draw[bpmn flow] (t) -- (p3);
  \draw[bpmn flow] (t) -- (p4);
  \node[font=\scriptsize\itshape, text width=3.6cm, align=center] at (1.5,-1.8) {both branches are always executed};
  \node[bpmn task, font=\scriptsize\sffamily] (r) at (7.2,0) {5.~rejection?};
  \node[bpmn task, font=\scriptsize\sffamily] (p6) at (10,0.7) {6.~reject letter};
  \node[bpmn task, font=\scriptsize\sffamily] (p7) at (10,-0.7) {7.~estimate};
  \draw[bpmn flow] (r) -- (p6);
  \draw[bpmn flow] (r) -- (p7);
  \node[font=\scriptsize\itshape, text width=3.6cm, align=center] at (8.7,-1.8) {exactly one branch is executed};
\end{tikzpicture}
\end{center}

Joins suffer the same ambiguity: \textsc{policy} and \textsc{premium} must \emph{both} complete before \textsc{rejection?} can fire (AND-join), whereas \textsc{filing} is reached by \textsc{reject letter} or by \textsc{payment} but never by both (XOR-join), and \textsc{estimate} is entered from \textsc{rejection?} or from \textsc{revision}: possibly executed, but not always both.
% NOTE: slide p.29 annotates the filing join as "tasks 6 and 13"; task 13
% does not exist: the converging branches are 6 (reject letter) and
% 11 (payment). Corrected here: suspected typo on the official slide.
The graph as drawn does not say which fork or join is which: the semantics lives in the analyst's head, not in the diagram. A formal notation has to remove this ambiguity.

% ── slides p.30--p.37 ───────────────────────────────────
\subsection{Disambiguation}

The fix is to type every split and join with a \emph{gateway}, a notation borrowed from BPMN: a diamond marked $+$ for parallel (AND) splits and joins, a diamond marked $\times$ for choice (XOR) splits and joins. The routing logic moves from the analyst's head into explicit nodes (one gateway per fork and join) and the arrows regain a uniform meaning:

\begin{center}
\resizebox{\linewidth}{!}{%
\begin{tikzpicture}
  \node[bpmn task, font=\scriptsize\sffamily] (rec) at (0,0) {1.~recording};
  \node[bpmn task, font=\scriptsize\sffamily] (typ) at (1.9,0) {2.~type};
  \node[bpmn gateway, font=\scriptsize] (a1) at (3.1,0) {$+$};
  \node[bpmn task, font=\scriptsize\sffamily] (pol) at (4.4,0.9) {3.~policy};
  \node[bpmn task, font=\scriptsize\sffamily] (pre) at (4.4,-0.9) {4.~premium};
  \node[bpmn gateway, font=\scriptsize] (a2) at (5.7,0) {$+$};
  \node[bpmn task, font=\scriptsize\sffamily] (rej) at (7.1,0) {5.~rejection?};
  \node[bpmn gateway, font=\scriptsize] (x1) at (8.5,0) {$\times$};
  \node[bpmn task, font=\scriptsize\sffamily] (let) at (10.3,1.2) {6.~reject letter};
  \node[bpmn gateway, font=\scriptsize] (x2) at (8.5,-1.4) {$\times$};
  \node[bpmn task, font=\scriptsize\sffamily] (est) at (10.3,-1.4) {7.~estimate};
  \node[bpmn task, font=\scriptsize\sffamily] (rea) at (10.3,-2.7) {8.~reaction};
  \node[bpmn task, font=\scriptsize\sffamily] (ass) at (12.5,-2.7) {9.~assessment?};
  \node[bpmn gateway, font=\scriptsize] (x3) at (14.2,-2.7) {$\times$};
  \node[bpmn task, font=\scriptsize\sffamily] (rev) at (12.5,-1.4) {10.~revision};
  \node[bpmn task, font=\scriptsize\sffamily] (pay) at (15.6,-2.7) {11.~payment};
  \node[bpmn gateway, font=\scriptsize] (x4) at (14.2,1.2) {$\times$};
  \node[bpmn task, font=\scriptsize\sffamily] (fil) at (15.8,1.2) {12.~filing};
  \draw[bpmn flow] (rec) -- (typ);
  \draw[bpmn flow] (typ) -- (a1);
  \draw[bpmn flow] (a1) -- (pol);
  \draw[bpmn flow] (a1) -- (pre);
  \draw[bpmn flow] (pol) -- (a2);
  \draw[bpmn flow] (pre) -- (a2);
  \draw[bpmn flow] (a2) -- (rej);
  \draw[bpmn flow] (rej) -- (x1);
  \draw[bpmn flow] (x1) -- (let);
  \draw[bpmn flow] (x1) -- (x2);
  \draw[bpmn flow] (x2) -- (est);
  \draw[bpmn flow] (rev) to[bend right=18] (x2);
  \draw[bpmn flow] (est) -- (rea);
  \draw[bpmn flow] (rea) -- (ass);
  \draw[bpmn flow] (ass) -- (x3);
  \draw[bpmn flow] (x3) -- (rev);
  \draw[bpmn flow] (x3) -- (pay);
  \draw[bpmn flow] (pay) -- (x4);
  \draw[bpmn flow] (let) -- (x4);
  \draw[bpmn flow] (x4) -- (fil);
\end{tikzpicture}}
\end{center}

Reading the gateways: the $+$ pair encloses the parallel checks; the first $\times$ routes to \textsc{reject letter} or \textsc{estimate}; a second $\times$ joins the loop back-edge from \textsc{revision} into \textsc{estimate}; a third splits after \textsc{assessment?} into repeat or exit; the last one merges the two alternative paths into \textsc{filing}.

% ── slides p.38--p.40 ───────────────────────────────────
Adding a start event and an end event, and enclosing the flow in a \emph{pool} labelled with the responsible organisation, turns the diagram into valid BPMN~2.0: drawable in any editor,\footnote{\url{https://camunda.com/download/modeler/}; \url{https://bpmn.io/}} exchangeable as \texttt{.bpmn} XML, executable by a compliant engine. A pool encloses a single participant of the choreography; inter-organisational interaction is modelled by adding more pools connected by message flows, a construct introduced in the next chapters. Tasks may further carry BPMN task-type markers.

% ── slide p.41 ──────────────────────────────────────────
The same control flow renders equally well in \textbf{EPC} (Event-driven Process Chain) syntax: tasks become rounded rectangles, events become hexagons, and the same $+$/$\times$ connectors type the forks and joins.\footnote{Drawn in ARIS Express: \url{https://ariscommunity.com/aris-express}}

% ── slide p.42 ──────────────────────────────────────────
In \textbf{workflow-net} form the encoding is structural: tasks become square transitions, branching points become explicit places, an XOR fork is several transitions competing for one place, an AND fork is one transition with several output places.\footnote{Drawn in WoPeD: \url{http://woped.dhbw-karlsruhe.de/}} The net has a single source place (\texttt{start}) and a single sink place (\texttt{end}). This is the formal counterpart that the Petri-net chapters analyse rigorously.

% ── slide p.43 ──────────────────────────────────────────
\subsection{Ingredients}

Any graphical process language must provide:

\begin{itemize}
  \item \textbf{nodes} for start and end: one source and one sink per process;
  \item \textbf{nodes} for tasks: the activities the process is made of;
  \item \textbf{nodes} for joins and splits, distinguishing concurrency (AND), internal decisions (XOR resolved by the system), and external decisions (XOR resolved by the environment: split only);
  \item \textbf{edges} for the links carrying causal and temporal dependencies.
\end{itemize}

Three further dimensions remain open: \emph{responsibility} (who owns the process and each task), \emph{information} (the data flowing along the edges), \emph{platform} (bindings to concrete services and runtimes), and the BPMN, EPC, and BPEL chapters fill them in.

% ── slide p.44 ──────────────────────────────────────────
Concrete notations are points in a design space: events as circles, stars, or hexagons; tasks as rectangles, ellipses, or gears; gateways as diamonds, rings, or bracket icons; links as solid, dashed, or double-headed arrows. No choice is better \emph{a priori}, as long as it is unambiguous and consistent.

% ── slides p.45--p.46 ───────────────────────────────────
\begin{examplebox}{Exercise: design a notation}
Invent your own diagrammatic notation to describe the following interaction protocol; choose symbols, shapes, and colours carefully. \textit{Alice wants to sell her car, Bob is interested in buying it. Alice asks some quote. Bob can accept the bargain, refuse it, or make a counteroffer. Alice can accept or make a counteroffer, and so on, until either the bargain is accepted or refused.}
\end{examplebox}

% ── slide p.47 ──────────────────────────────────────────
The winning entry of the 2024/25 edition of this exercise was a two-lane sequence diagram: one vertical lifeline per participant, the message types (\textsc{makes an offer}, \textsc{counteroffer}, \textsc{accepts}, \textsc{refuses}) drawn as arrows between the lifelines and colour-coded by originator, and a shared \textsc{end} node closing the protocol. The bargaining loop emerges from the arrow topology alone, with no explicit gateway.

% ── slides p.48--p.49 ───────────────────────────────────
% SKIP: section_break (p.48), duplicate of p.44 (p.49)

% ── slide p.50 ──────────────────────────────────────────
\subsection{A sneak peek at standards}

The notations studied in this course (BPMN, EPC and workflow nets) agree on the ingredient list but disagree on the symbols:

\begin{center}
\begin{tabular}{lccc}
\toprule
 & \textbf{BPMN} & \textbf{EPC} & \textbf{WfN} \\
\midrule
start / end & circle & hexagon & circle (place) \\
join / split & diamond & circle & square (transition) \\
tasks       & rounded rect. & rounded rect. & square (transition) \\
links       & solid arrow & dashed arrow & solid arrow \\
\bottomrule
\end{tabular}
\end{center}

The course's notational baseline is BPMN; EPC and workflow nets enter the picture in the analysis chapters.

% ── slides p.51--p.52 ───────────────────────────────────
\subsection{Alice--Bob car selling in BPMN}

As a BPMN collaboration, each participant gets a pool, and the offer/counteroffer/accept/refuse exchange travels on message flows between the pools. Alice's pool is sketched below; Bob's mirrors it, and the message flows between the two are omitted:

\begin{center}
\begin{tikzpicture}
  \node[bpmn event] (st) at (0,0) {};
  \node[bpmn task, font=\scriptsize\sffamily] (ask) at (1.6,0) {Ask quote};
  \node[bpmn gateway, font=\scriptsize] (xm) at (3.3,0) {$\times$};
  \node[bpmn gateway] (eb) at (4.6,0) {};
  \node[bpmn event, minimum size=0.5em] at (eb.center) {};
  \node[bpmn task, font=\scriptsize\sffamily, text width=1.6cm, align=center] (rc) at (6.7,0) {Receive counteroffer};
  \node[bpmn task, font=\scriptsize\sffamily, text width=1.6cm, align=center] (ev) at (9.3,0) {Evaluate proposal};
  \node[bpmn gateway, font=\scriptsize] (xs) at (11.4,0) {$\times$};
  \node[bpmn event, line width=1.1pt, label={[font=\scriptsize]below:accept}] (acc) at (13.2,0) {};
  \node[bpmn task, font=\scriptsize\sffamily, text width=1.6cm, align=center] (sc) at (11.4,-1.6) {Send counteroffer};
  \node[bpmn event, label={[font=\scriptsize]below:deal done}] (dd) at (5.6,-2.8) {};
  \node[bpmn event, line width=1.1pt] (e1) at (7.6,-2.8) {};
  \node[bpmn event, label={[font=\scriptsize]below:bargain refusal}] (rf) at (9.6,-2.8) {};
  \node[bpmn event, line width=1.1pt] (e2) at (11.6,-2.8) {};
  \draw[bpmn flow] (st) -- (ask);
  \draw[bpmn flow] (ask) -- (xm);
  \draw[bpmn flow] (xm) -- (eb);
  \draw[bpmn flow] (eb) -- (rc);
  \draw[bpmn flow] (rc) -- (ev);
  \draw[bpmn flow] (ev) -- (xs);
  \draw[bpmn flow] (xs) -- (acc);
  \draw[bpmn flow] (xs) -- (sc);
  \draw[bpmn flow] (sc.west) -| (xm);
  \draw[bpmn flow] (eb) |- (dd);
  \draw[bpmn flow] (dd) -- (e1);
  \draw[bpmn flow] (eb.south) to[bend right=25] (rf);
  \draw[bpmn flow] (rf) -- (e2);
  \begin{scope}[on background layer]
    \node[draw=accent, thick, inner sep=10pt, fit=(st)(acc)(sc)(dd)(e2)] (pool) {};
  \end{scope}
  \node[font=\scriptsize\bfseries\sffamily, anchor=north west] at (pool.north west) {Alice};
\end{tikzpicture}
\end{center}

After \textsc{ask quote}, Alice does not decide internally what happens next: she waits for whichever message arrives first: a counteroffer to evaluate, an acceptance, a refusal. The diamond with the inner circle is BPMN's \emph{event-based gateway}: the routing decision is taken by the environment, not by the process. If Alice evaluates a counteroffer, she either accepts (closing the deal) or sends a counteroffer of her own and waits again. The event-based gateway realises the \emph{deferred choice}: the external-decision ingredient listed above; internal and external choice return rigorously with Petri nets in the next chapter.
